\documentclass[11pt]{article}
\usepackage[review]{acl}
\usepackage{times}
\usepackage{latexsym}
\usepackage[T1]{fontenc}
\usepackage[utf8]{inputenc}
\usepackage{microtype}
\usepackage{inconsolata}
\usepackage{lipsum}

\title{Cross-Lingual Parsing with Non-ASCII Author Names}

\author{First Author \\
  Test University \\
  \texttt{first@test.edu} \\\And
  Second Author \\
  Test University \\
  \texttt{second@test.edu} \\}

\begin{document}
\maketitle

\begin{abstract}
This paper tests reference extraction when author names contain diacritics,
umlauts, and other non-ASCII Unicode characters that appear frequently in
NLP author names. Incorrect transliteration of these characters causes
false-negative author matching during hallucination detection.
\end{abstract}

\section{Introduction}
\lipsum[1-2]
Multilingual NLP has benefited greatly from cross-lingual models
\cite{schuetze-1992-dimensions,sogaard-2011-semisupervised,nivre-etal-2016-universal}.
Character-level representations \cite{bojanowski-etal-2017-enriching} handle
morphologically rich languages effectively.

\lipsum[3-4]

\section{Background}
\lipsum[5-7]
Transfer learning across languages \cite{conneau-etal-2020-unsupervised,wu-dredze-2019-beto}
has become the dominant paradigm.

\lipsum[8-9]

\section{Methods}
\lipsum[10-12]
We evaluate on Universal Dependencies \cite{nivre-etal-2016-universal}
following the methodology of \cite{sogaard-2011-semisupervised}.

\lipsum[13-14]

\section{Conclusion}
\lipsum[15]

\bibliography{unicode_refs}
\bibliographystyle{acl_natbib}

\end{document}
